For $a\in A$, the monic polynomial $\prod_{\sigma\in G}(T-\sigma(a))$ has coefficients fixed by $G$ and annihilates $a$. Thus $A$ is integral over $A^G$.

The formula $\sigma(a/s)=\sigma(a)/\sigma(s)$ gives an action on $S^{-1}A$. The natural map $(S\cap A^G)^{-1}A^G\to(S^{-1}A)^G$ is injective: if $a/s$ maps to zero, some $t\in S$ annihilates $a$, and the invariant element $\prod_{\sigma\in G}\sigma(t)\in S\cap A^G$ also annihilates it.

Given an invariant fraction, multiply numerator and denominator by the conjugates of its denominator to write it as $a/d$ with $d\in S\cap A^G$. For each $\sigma\in G$, choose $t_\sigma\in S$ with $t_\sigma(\sigma(a)-a)=0$. Put $c=\prod_{\tau\in G}\tau(\prod_{\sigma\in G}t_\sigma)$. Then $c\in S\cap A^G$ annihilates every $\sigma(a)-a$, so $ca\in A^G$. Thus $a/d=ca/(cd)$ comes from the invariant localization.
